%!TEX program = uplatex
\RequirePackage{plautopatch}
\documentclass[uplatex,dvipdfmx]{jlreq}
\usepackage{amsmath,amssymb,amscd,amsthm}
\usepackage{enumerate}
\usepackage{url}
\newcommand{\RR}{\mathop {\Bbb R}\nolimits}
\newcommand{\FF}{\mathop {\Bbb F}\nolimits}
\newcommand{\Id}{\mathop {\rm Id}\nolimits}
\newcommand{\ZZ}{\mathop {\Bbb Z}\nolimits}
\newcommand{\PP}{\mathop {\Bbb P}\nolimits}
\newcommand{\QQ}{\mathop {\Bbb Q}\nolimits}
\newcommand{\CC}{\mathop {\Bbb C}\nolimits}
\newcommand{\NN}{\mathop {\Bbb N}\nolimits}
\newcommand{\Gal}{\mathop {\textrm Gal}\nolimits}

\pagestyle{empty}

\begin{document}

%======================================================================
% 講演１：難波誠「分岐被覆のいろいろな例」
%======================================================================
\begin{center}
{\Large\bfseries 分岐被覆のいろいろな例}

\hspace*{\fill}{\large\bfseries 難波 誠（阪大・理）}
\end{center}

\bigskip

分岐被覆空間 $X$ とは、
あたえられた空間（底空間）$M$ 上に分岐しつつ拡がっている
空間のことであり、分岐被覆写像とは、
その様子を表す $X$ から $M$ への写像のことである。
ここで底空間は良く知られた空間をとるのが普通である。
分岐被覆の研究はいわば地上から天上のきらめく星を見上げて観察し研究するようなものである。

（分岐の無い）被覆の理論は通常、位相幾何学の一章として論じられるが、その一
般化である分岐被覆の理論は大抵省略される。ところが実際の例として現れ調べられ
るのは、分岐被覆の方が多い。例えば、有理関数は複素数球面から
それ自身への有限分岐被覆写像とみなしうる。代数関数は閉リーマン面から複素数球
面への有限分岐被覆写像とみなしうる。（歴史的にはリーマンが逆に、複素数球面上
の分岐被覆空間としてリーマン面を捉えた。）
楕円関数、モジュラー関数もそれぞれ複素数平面、上半平面から複素数球面への、無
限分岐被覆写像とみなしうる。

ガロア理論を、体の拡大を群で統制する理論とみるとき、代数多様体上の分岐被覆
の理論は、（多変数）代数関数体のガロア理論の幾何学と言える。
ガロア自身がその遺書で「あいまいの理論」と呼んだ謎の理論はその予見だったのではないかと言われている。

このように分岐被覆は代数幾何と密接に結びついているが、一方、局所化を考える
ことにより、特異点論とも結びつく。他方、（ベリーの定理など）いろいろな意味で
数論にも関係する。群論、表現論、フックス型微分方程式、複素関数論にも関係する。
また分岐被覆のモジュライを考えることにより、不変式論にも関係する。もちろん分
岐被覆の基礎にトポロジーがあることは言うまでもない。

分岐被覆はこのように、いろいろな分野の交叉する地点にあり、逆にこれらの分野
に重要な実例をあたえている。

この講演は、トップバッターとして、分岐被覆の基礎理論を述べると同時に、この
地点にさまざまな分野が交叉している様子を実例を主体に紹介することを目的とする。

（分岐の無い）被覆が底空間の基本群で統制されるように、分岐被覆は分岐点の集
合の補空間の基本群で統制される。そのため、分岐被覆の立場からも、基本群を調べ
ることは重要である。基本群は底空間が複素１次元ならば簡単に
記述されるが、複素２次元以上になると、急に難しくなり計算も大変になる。ホモロ
ジーやコホモロジーなど、可換的なものよりはるかに難しい。基本群はしかし、不思
議な神秘的な対象で、我々の興味をそそり、その探究へとかりたてる。

基本群にかんしては、この講演をひきつぎ岡睦雄氏が論じる予定である。

\begin{center}\textsc{参考文献}\end{center}

\begin{enumerate}[{[1]}]
  \setlength{\itemsep}{2pt}
  \item[{[1]}] M.~Namba,
        \textit{Branched Coverings and Algebraic Functions},
        Pitman Research Notes in Math.\ 161, Pitman--Longman, 1987.
  \item[{[2]}] M.~Namba,
        Finite branched coverings of complex manifolds,
        \textit{Sugaku Expositions} \textbf{5} (1992), 193--211.
\end{enumerate}

\newpage

%======================================================================
% 講演２：岡睦雄「基本群入門」
%======================================================================
\begin{center}
{\Large\bfseries 基本群入門}

\hspace*{\fill}{\large\bfseries 岡 睦雄（都立大・理）}
\end{center}

\bigskip

基本群は位相空間 $X$ の非可換な不変量として一番大事なものである。
それ自体が研究対象として面白いだけでなく、
空間 $X$ の上の（分岐）被覆写像やファイバー束の構造などの研究に欠かせない。
$n$ 次元代数多様体は $n$ 次元射影空間 $P(n)$ 上の有限被覆として捉えられる。
その際分岐集合として超曲面 $V$ が現れる。
故に $P(n)\setminus V$ の基本群の研究が重要となってくる。
ザリスキーはこの基本群が一般の超平面で切断することによって、
２次元の射影空間の中の射影曲線 $C$ の補空間の基本群に等しいことを示して、
$P(2)\setminus C$ の基本群の組織的な研究をした。

この講義では $P(2)\setminus C$ の基本群に焦点を絞って、
その基本的な性質（１次元ホモロジーとの関係、ミルナー束との関係）、
計算方法（van-Kampen、Zariski のペンシル method）、
局所特異点との関係（A.~Libgober の定理、R.~Randell の定理）、
基本群から得られる不変量としての Alexander 多項式（Esnault、Libgober、Artal、
Loeser--Vaquié）などを解説する。
特に基本群が可換なのはいつか、非可換なら特異点はどうなるか？などの初等的な解説をする。
時間の余裕があれば応用についても少し話す。

\newpage

%======================================================================
% 講演３：難波誠「閉リーマン面の自己同型群と分岐被覆」
%======================================================================
\begin{center}
{\Large\bfseries 閉リーマン面の自己同型群と分岐被覆}

\hspace*{\fill}{\large\bfseries 難波 誠（阪大・理）}
\end{center}

\bigskip

示性数が２以上の閉リーマン面の自己同型群が有限群になることはよく知られている。
しかるにグリーンバーグは1963年につぎの定理をアナウンスしている：
「定理：任意の非自明有限群 $G$ に対し、その自己同型群が $G$ と同型になる閉リーマン面が存在する。」
しかしその証明の詳細は述べられていない。この講演では、複素数球面上のガロア
分岐被覆の性質、フックス群の性質などを援用した証明をあたえる。この証明は複素
数球面の有限ガロア分岐被覆のモジュライ空間に新しい知見を与えると考える。

\newpage

%======================================================================
% 講演４：島田伊知朗「開代数多様体の基本群」
%======================================================================
\begin{center}
{\Large\bfseries 開代数多様体の基本群}

\hspace*{\fill}{\large\bfseries 島田 伊知朗（北大・理）}
\end{center}

\bigskip

Zariski--van Kampen の定理は開代数多様体の基本群を調べる上で
もっとも大切な道具のひとつである。
本講演では、この定理のさまざまなヴァリアントをいくつかの応用とともに紹介する。

\begin{center}\textsc{参考文献}\end{center}

\begin{enumerate}[{[1]}]
  \setlength{\itemsep}{2pt}
  \item[{[1]}] I.~Shimada,
        On Zariski--van Kampen theorem,
        preprint, \url{http://www.math.sci.hokudai.ac.jp/~shimada/preprints.html}.
  \item[{[2]}] I.~Shimada,
        The fundamental group of the complement of a resultant hypersurface,
        preprint, \url{http://www.math.sci.hokudai.ac.jp/~shimada/preprints.html}.
  \item[{[3]}] I.~Shimada,
        Fundamental groups of algebraic fiber spaces,
        preprint, \url{http://www.math.sci.hokudai.ac.jp/~shimada/preprints.html}.
\end{enumerate}

\newpage

%======================================================================
% 講演５：徳永浩雄「3次方程式の幾何学」
%======================================================================
\begin{center}
{\Large\bfseries 3次方程式の幾何学 ---$S_3$-被覆と平面$6$次曲線の補空間の基本群---}

\hspace*{\fill}{\large\bfseries 徳永 浩雄（都立大・理）}
\end{center}

\bigskip

$3$ 次方程式
$x^3 + ax + b = 0$ の係数
$a, b$ が $4a^3 + 27b^2 = 0$
をみたすとき、この方程式が重根をもつという事実はよく知られている。
係数 $a, b$ が $(a, b)$-平面を動くパラメータであると考えれば、この事実は $(a,b)$-平面を $3$ 重に覆う被覆が
曲線 $4a^3 + 27b^2 = 0$ の沿って分岐しているということに他ならない。
$(a, b)$-平面を射影平面 $\PP^2$ に、
曲線 $4a^3 + 27 b^2 = 0$ を一般の平面曲線 $B$ に置き換え、
$B$ で分岐する Galois でない $3$ 次被覆、言い換えれば、
$B$ で分岐する $3$ 次対称群を Galois 群とする分岐被覆、が存在するか
という問題に関して講演する。

\medskip

基本群の言葉に翻訳すればこの問題は、
「$B$ の補空間 $\PP^2\setminus B$ の基本群から $3$ 次対称群
$S_3$ への全射が存在するか否かを決定せよ」となる。
が、$\pi_1(\PP^2\setminus B)$ を求めることは
思いの外、難しくまた微妙である。例えば Zariski の1929年の論文にある例
---$6$ 個のカスプをもつ $6$ 次曲線--- では
基本群はその特異点の位置によって変わってしまう。

\medskip

本講演では $3$ 次被覆を $3$ 次方程式の解法にこだわった観点から捉えてみる。
$3$ 次方程式に解の公式があるということは良く知られている。
その「精神」を $3$ 次被覆の構成に生かし、
そうした手法で解説できる現象や事実---例えば以下の\S2, 3にあるような---を
解説したい。$3$ 次方程式といってもなかなか侮れないなと感じてもらえれば幸いである。

\medskip

\S1\quad $S_3$-被覆の攻略法

\medskip

\S2\quad $S_3$-被覆の存在に関するあるひとつの条件

\medskip

\S3\quad 楕円曲面の Mordell--Weil 群の位数 $3$ の元と $S_3$-被覆：Zariski 対の例

\begin{center}\textsc{参考文献}\end{center}

\begin{enumerate}[{[1]}]
  \setlength{\itemsep}{2pt}
  \item[{[1]}] 徳永浩雄・島田伊知朗 他，
        「代数曲線と特異点」，特異点の数理 4，共立出版，2001年．
\end{enumerate}

\end{document}
